\documentclass[11pt,fleqn]{article}
\usepackage[margin=0.8in]{geometry}
\usepackage{amsmath,amssymb}
\usepackage[hidelinks]{hyperref}
\setlength{\mathindent}{0pt}

\newcommand{\X}{\mathcal X}
\newcommand{\C}{\mathcal C}
\newcommand{\Hy}{\mathcal H}
\newcommand{\err}{\operatorname{error}}
\newcommand{\eps}{\epsilon}
\newcommand{\Adv}{\operatorname{Adv}}

\begin{document}
\raggedright
\setlength{\parskip}{0.6em}

{\LARGE\bfseries Page 6: Learning Parity and Interpreting the Demos\par}
\medskip

\section{Sample Bound Calculator}

Play around with different values of the parameters.

\section{Oracle Challenge}

Try playing around with this. In the keyed branch, the browser will generate a secret HMAC key for the round and then compute HMAC-SHA-256 for each input, returning the first bit of result~\cite{mdn}.

\section{A Family We CAN Learn Efficiently}

A parity function is one that has a secret vector $s \in \{0, 1\}^n$~\cite{mit}. It labels an input $x$ by computing
\[f_s(x) = \sum_{i=1}^n s_i x_i \pmod 2\]

Notice how a single example $(x, y)$ gives an equation $x \cdot s = y$. If we have multiple examples, this yields a system of equations, representable in matrix form as $Xs = y$, where each row of $X$ is an input and each element of $y$ is its corresponding label. Note that this is a system over GF(2). Thus, once we get a matrix with rank $n$, we can recover the unique $s$ in polynomial time using Gaussian elimination.

\begin{thebibliography}{9}
\bibitem{mit}
MIT 6.S976 / 18.S996.
\emph{Cryptography and Machine Learning: Foundations and Frontiers}.
\href{https://mlcrypto.mit.edu/course/}{Lectures 4--5}, Spring 2026.

\bibitem{mdn}
MDN Web Docs.
\href{https://developer.mozilla.org/en-US/docs/Web/API/Web_Crypto_API}
{Web Crypto API} and
\href{https://developer.mozilla.org/en-US/docs/Web/API/SubtleCrypto/sign#hmac}
{\texttt{SubtleCrypto.sign()}, HMAC}.
\end{thebibliography}

\end{document}
